\section{Introduction}

\paragraph{Background.} A mobile GUI agent executes a task by repeatedly looking at a screen and emitting an action. On a long cross-app task the screens it has already seen scroll out of view, and some of them carried information it will need later. An agent that reads a recipe page, then opens a shopping list, must type ``vanilla extract'' nine steps after the only screen that ever showed those two words. We call a step whose required text appeared on an earlier screen, and only there, a \emph{delayed information need}; the screen that held the text is the \emph{source screen}, and the number of screens since is the \emph{gap}. Section~\ref{sec:needs} shows these steps are common in long tasks: 31\% of sampled GUI-Odyssey~\citep{guiodyssey} trajectories contain at least one, with a median gap of 5 screens. They are absent from the short-task benchmarks the field develops on: AITZ~\citep{aitz} has none in 18,643 steps, consistent with tasks that average 7.4 steps.

\paragraph{Existing approaches and their limitation.} Two designs recover such text, and both are prompted rather than trained. \emph{Full history} keeps the OCR text of every earlier screen in context and asks the model to name the source screen and the text in one call. \emph{Gist-and-lookup}, introduced for long documents by ReadAgent~\citep{readagent} and instantiated here as a \emph{screen index}, writes one short line per screen and asks a reader to pick a screen from the lines before re-reading it in full. Both degrade as the history grows: full history localizes the right screen in 68\% of needs that arise by step 15 and in 46\% of later ones, and the screen index reaches only 41\% overall. The shared limitation is not the representation but the supervision. Neither design has been trained for this decision, because training it means labeling which earlier screen an agent should re-read, and the one trained look-back agent we know of, PAL-UI~\citep{palui}, pays for that with 8.6K curated step-level instructions. On a new corpus, that cost is paid again.

\paragraph{The challenge.} Labeling look-back by hand is expensive for a structural reason: the label is not a property of one screen but of a relation between a step and all the screens before it, so a human annotator must read an entire trajectory to label one decision. A cheap automatic label would have to identify, for a given step, exactly which of the $t$ preceding screens could answer it. It must also be trustworthy enough to train on, which rules out heuristics such as recency or lexical overlap, both of which we show are wrong often enough to mislead a learner (Section~\ref{sec:index}).

\paragraph{Our method.} Such a label already exists, as a by-product of defining the problem. A delayed information need is verified by a string check: the typed text appears in the OCR of an earlier screen and not on the current one. Run that same check against every earlier screen rather than stopping at the first hit, and it returns the set of screens that contain the text. That set is a complete, exact label over all $t$ candidates, obtained at the cost of one substring test per screen, with no human in the loop. We call training a look-back policy on these labels \emph{verifier-supervised look-back}, and we train two instantiations on them (Figure~\ref{fig:method}): a \emph{per-screen scorer} that reads one earlier screen at a time and outputs the probability that it holds the needed text, and a \emph{full-history reader} that sees the whole history in one call and emits a screen number. A second call, the \emph{read stage}, then types the text.

\paragraph{Contributions.}
\begin{itemize}\setlength{\itemsep}{1pt}
\item \textbf{Problem and free labels.} We define delayed information needs, give an automatic verifier that both mines and labels them, and measure their prevalence: absent from AITZ, present in 31\% of sampled GUI-Odyssey trajectories (Section~\ref{sec:needs}).
\item \textbf{Method.} Verifier-supervised look-back, with two instantiations that share one label set and one objective and differ only in how the history reaches the model (Section~\ref{sec:method}).
\item \textbf{Finding: supervision, not architecture, is the lever.} The labels lift the full-history reader by 10.6 localization points and the per-screen scorer by 22.3, and reverse which architecture is better; a paired test separates the two trained variants on neither metric (Section~\ref{sec:main}). We locate the decomposition's real gain at gaps of 8--12 screens, where it is worth 21.8 points (Section~\ref{sec:gap}).
\item \textbf{External validation and honest limits.} On a fourth sample held out from every design decision, both trained variants keep a 17--20 point margin over prompting and remain indistinguishable (Section~\ref{sec:ext}). Three ablations show the design is insensitive to its prompt and loss settings (Section~\ref{sec:ablation}). At gaps of 13 or more screens both trained variants fall \emph{below} prompting, which we report rather than omit (Section~\ref{sec:limits}).
\item \textbf{Two mandatory evaluation filters.} Scoring only the earliest source screen biases one arm by 15 points, and admitting cases whose answer is also on the current screen reverses the sign of a 17-point result. We quantify both and give the filter conditions (Section~\ref{sec:index}).
\end{itemize}
